\chapter{Blood Glucose Test}
\index{Blood Glucose Test}

\section*{Definition and Overview}
A blood glucose test measures the concentration of glucose (sugar) in the blood. Glucose is the body's principal source of energy, and its level in the bloodstream reflects the balance between dietary intake, glucose production by the liver, and uptake by tissues under the influence of insulin and other hormones. Blood glucose testing is used to screen for and diagnose diabetes mellitus and prediabetes, to monitor glucose control in people already living with diabetes, and to investigate symptoms such as unexplained thirst, fatigue, or weight loss. Several distinct forms of testing exist: the fasting plasma glucose test, the random (or casual) plasma glucose test, the oral glucose tolerance test (OGTT), and the glycated haemoglobin (HbA1c) test, which estimates average blood glucose over the preceding two to three months. In addition, home glucose monitors (finger-prick devices and, increasingly, continuous glucose monitoring sensors) allow people with diabetes to measure their own glucose levels between clinic visits, supporting day-to-day treatment decisions.

\section*{Epidemiology}
Diabetes is one of the most common chronic diseases worldwide and is increasing in prevalence. The World Health Organization estimates that roughly one in eleven adults lives with diabetes, and a substantial proportion of cases remain undiagnosed. Type 2 diabetes is by far the most common form, and its frequency rises with age, obesity, physical inactivity, and family history, while type 1 diabetes typically begins in childhood or young adulthood. Prediabetes -- a state in which glucose levels are above normal but below the diagnostic threshold for diabetes -- is also common and represents a critical window for prevention. Blood glucose testing underpins both the detection and the ongoing management of these conditions, and millions of people worldwide perform home glucose monitoring every day. Because many cases are silent, guideline-based screening of at-risk adults is an important public health measure.

\section*{Aetiology and Pathophysiology}
Glucose is derived from carbohydrate digestion and from the liver, which releases stored glucose between meals. Insulin, secreted by the beta cells of the pancreatic islets, drives glucose into muscle and fat cells and suppresses liver glucose output. In diabetes this regulatory system fails: in type 1 diabetes the immune system destroys the insulin-producing beta cells, leading to an absolute insulin deficiency, whereas in type 2 diabetes tissues become resistant to insulin and the pancreas eventually cannot produce enough to compensate. Blood glucose tests therefore measure how effectively this system maintains glucose homeostasis.

\begin{itemize}
    \item \textbf{Fasting plasma glucose:} measured after at least 8 hours without food; reflects basal glucose control
    \item \textbf{Random plasma glucose:} measured at any time regardless of food intake; useful in the symptomatic patient
    \item \textbf{Oral glucose tolerance test:} a fasting sample, a glucose drink, then further samples (usually at 1 and 2 hours) to assess glucose disposal
    \item \textbf{HbA1c:} reflects average glucose over 2--3 months, does not require fasting, and can be measured at any time of day
    \item \textbf{Home monitoring:} capillary finger-prick testing or sensor-based continuous glucose monitoring performed by the person with diabetes
\end{itemize}

A single high reading may reflect stress, acute illness, trauma, or medicines such as corticosteroids, rather than diabetes, whereas HbA1c smooths out day-to-day variation and gives a more stable picture of long-term control. Conversely, conditions that alter red blood cell survival, such as anaemia, haemoglobinopathies, or recent blood loss, can render HbA1c unreliable.

\section*{Clinical Features}
Most blood glucose tests involve a simple blood sample with only brief discomfort, but the clinical context in which testing occurs is important.

\begin{itemize}
    \item Blood taking causes a brief stinging sensation and may leave a small bruise
    \item Fasting tests require avoiding food (and sometimes adjusting medicines) beforehand
    \item The oral glucose tolerance test takes about two hours, and the glucose drink is very sweet
    \item Symptoms that prompt testing include excessive thirst, frequent urination, fatigue, unintentional weight loss, and blurred vision
    \item Hypoglycaemia (a glucose level too low) can cause sweating, trembling, palpitations, hunger, confusion, and, if severe, seizures or loss of consciousness
\end{itemize}

People with type 1 diabetes may present acutely with weight loss, abdominal pain, vomiting, and a fruity or chemical smell on the breath, reflecting diabetic ketoacidosis. In contrast, type 2 diabetes is often asymptomatic and discovered incidentally during routine testing.

\section*{Differential Diagnosis}
Interpreting a glucose result requires care, because an elevated reading does not always mean diabetes:

\begin{itemize}
    \item \textbf{Prediabetes:} impaired fasting glucose or impaired glucose tolerance, often coexisting with obesity and metabolic syndrome
    \item \textbf{Stress hyperglycaemia:} transient elevation during acute illness, surgery, trauma, or myocardial infarction
    \item \textbf{Steroid-induced hyperglycaemia:} from glucocorticoid therapy
    \item \textbf{Gestational diabetes:} glucose intolerance first recognised during pregnancy
    \item \textbf{Endocrine causes:} Cushing's syndrome, phaeochromocytoma, acromegaly, or hyperthyroidism
    \item \textbf{Pancreatic disease:} chronic pancreatitis, cystic fibrosis, haemochromatosis, or pancreatic surgery
    \item \textbf{Preanalytical error:} non-fasting sampling, testing after a meal, or haemolysed samples
    \item \textbf{Conditions affecting red cell survival:} anaemia or haemoglobin variants can distort HbA1c results
\end{itemize}

\section*{When to Seek Medical Care}
See a doctor for testing if you have symptoms of high blood sugar -- increased thirst, frequent urination, unexplained weight loss, fatigue, or blurred vision -- or if you have risk factors such as a family history of diabetes, obesity, high blood pressure, or a history of gestational diabetes. Anyone with known diabetes who has repeated unexplained high readings, or who is taking glucose-lowering medicines and experiencing frequent low episodes, should also seek review.

\textbf{Contact your local emergency services for:}
\begin{itemize}
    \item Confusion, unusual behaviour, or loss of consciousness (possible severe hypoglycaemia or hyperglycaemic emergency)
    \item Difficulty breathing, especially with a fruity or chemical smell on the breath (possible ketoacidosis)
    \item Persistent vomiting with abdominal pain in a person with diabetes
    \item Seizures or an unresponsive person with known diabetes
\end{itemize}

\section*{Diagnosis and Investigations}
Diagnosis begins with a careful history -- symptoms, risk factors, medications, and family history -- followed by confirmatory laboratory testing, since an abnormal result should generally be confirmed on a second occasion in asymptomatic people.

\begin{itemize}
    \item \textbf{Fasting plasma glucose:} a fasting level at or above the diagnostic threshold indicates diabetes
    \item \textbf{Oral glucose tolerance test:} used when fasting results are borderline or during pregnancy
    \item \textbf{HbA1c:} an elevated result, confirmed by repeat testing, establishes the diagnosis without fasting
    \item \textbf{Urine tests:} glucose and ketones may be checked, particularly at presentation
    \item \textbf{Repeat testing:} a single abnormal result is usually confirmed with a second test
    \item \textbf{Related checks:} kidney function, lipid profile, blood pressure, and eye and foot examinations once diabetes is diagnosed
\end{itemize}

\section*{Management}
A blood glucose test is a diagnostic and monitoring tool; its value lies in how the result guides management:

\begin{itemize}
    \item \textbf{Normal results:} reassuring, but maintain healthy habits and repeat screening as advised
    \item \textbf{Prediabetes:} modest weight loss, dietary changes, and regular physical activity can prevent progression to diabetes
    \item \textbf{Diabetes:} lifestyle modification combined with glucose-lowering medicines or insulin as required, and regular monitoring
    \item \textbf{Home monitoring:} the frequency and method depend on the type of diabetes and the treatment regimen
    \item \textbf{Hypoglycaemia:} treat promptly with fast-acting carbohydrate and confirm recovery with a repeat reading
    \item \textbf{Follow-up:} regular HbA1c checks, often every 3--6 months, to assess long-term control
    \item \textbf{Education:} learn how to prepare for tests, interpret results, and adjust treatment safely with your healthcare team
\end{itemize}

\section*{Complications}
\begin{itemize}
    \item Bruising or, rarely, infection at the blood-taking site
    \item Hypoglycaemia from excessive glucose-lowering treatment, missed meals, or fasting
    \item Missed diagnosis if an abnormal result is not acted upon
    \item False reassurance from a single normal result in a person at high risk
    \item Prolonged poor control leads to microvascular damage (eyes, kidneys, nerves) and macrovascular disease (heart attack, stroke)
\end{itemize}

\section*{Prognosis and Outcomes}
Testing is quick, safe, and well tolerated, and the results are highly actionable. When diabetes is detected early and treated effectively, the risk of complications is greatly reduced, and many people with diabetes live long, healthy lives. HbA1c values that remain within the agreed target range are associated with fewer diabetes-related complications. Prediabetes can frequently be reversed with lifestyle change, and early detection of type 2 diabetes allows timely intervention. The overall benefit of any glucose test depends on acting on the result together with your healthcare team.

\section*{Prevention and Self-Care}
\begin{itemize}
    \item Have your risk assessed and glucose tested if you have risk factors or symptoms
    \item Follow instructions for fasting and for holding medicines before a scheduled test
    \item Maintain a healthy weight and be physically active to prevent type 2 diabetes
    \item Eat a balanced diet low in added sugar and refined carbohydrates
    \item If you have diabetes, attend regular HbA1c checks and eye, kidney, and foot reviews
    \item Learn to recognise and treat low blood sugar if you take glucose-lowering medicines
    \item Keep a record of home readings and bring it to appointments
\end{itemize}

\section*{Sources and Further Reading}
The following sources provide reliable, up-to-date information on blood glucose testing and diabetes:

\begin{itemize}
    \item NHS. \textit{NHS guidance on blood sugar testing and diabetes.} https://www.nhs.uk/conditions/
    \item Mayo Clinic. \textit{Information on diabetes diagnosis and blood glucose tests.} https://www.mayoclinic.org/
    \item World Health Organization. \textit{Diabetes fact sheets and diagnostic guidance.} https://www.who.int/
    \item Centers for Disease Control and Prevention. \textit{Diabetes information and statistics.} https://www.cdc.gov/
    \item MedlinePlus. \textit{Blood glucose test information for patients.} https://medlineplus.gov/
    \item NICE. \textit{Type 2 diabetes prevention and management guidelines.} https://www.nice.org.uk/
\end{itemize}
